\documentclass[conference]{IEEEtran}
\IEEEoverridecommandlockouts

\usepackage{cite}
\usepackage{amsmath,amssymb,amsfonts}
\usepackage{algorithm}
\usepackage{algorithmic}
\usepackage{graphicx}
\usepackage{textcomp}
\usepackage{xcolor}
\usepackage{booktabs}
\usepackage{multirow}
\usepackage{url}
\usepackage{hyperref}

\hypersetup{
  colorlinks=true,
  linkcolor=blue,
  citecolor=blue,
  urlcolor=blue,
  pdftitle={moe-engine: A Fault-Tolerant Runtime for Hyperscale Mixture-of-Experts Training},
  pdfauthor={Min Htet Myet}
}

\def\BibTeX{{\rm B\kern-.05em{\sc i\kern-.025em b}\kern-.08em
    T\kern-.1667em\lower.7ex\hbox{E}\kern-.125emX}}

\begin{document}

\title{moe-engine: A Fault-Tolerant Runtime for\\Hyperscale Mixture-of-Experts Training}

\author{
\IEEEauthorblockN{Min Htet Myet}
\IEEEauthorblockA{
\textit{Independent Researcher}\\
\texttt{github.com/Mattral/Composed-Mixture-of-Experts-Engine}
}
}

\maketitle

\begin{abstract}
We present moe-engine, a research-grade infrastructure runtime for training
large Mixture-of-Experts (MoE) models under continuous node failure at
hyperscale. The system combines a fused Triton routing kernel, composable
4D parallelism (data, expert, tensor, pipeline), invariant-checked
correctness, and elastic fault tolerance with automatic expert resharding.
This paper reports two developments beyond our prior work: first, the
system's first sustained validation on real GPU hardware (NVIDIA T4),
which surfaced and corrected a kernel compilation defect invisible to
CPU-only continuous integration and yielded measured throughput results
including an 80.1$\times$ speedup over a CPU reference path and a 100\%
pass rate on a storage-fault chaos scenario; second, a structural
refactoring replacing a 1{,}165-line monolithic module with eight
single-responsibility components, a migration to statically validated
configuration, and a Switch-Transformer-style expert capacity enforcement
mechanism. We additionally report, as a methodological contribution, a
root-caused case study of a configuration validator that silently degraded
to a no-op across an entire release cycle due to an undeclared dependency
--- a failure mode we argue generalizes to any system that pairs
safety-critical validation with an optional-dependency fallback. The test
suite grew from 148 to 348 passing cases over this arc, incorporating
property-based testing and a mocked collective-communication backend. We
report multi-node GPU benchmarking and full chaos-recovery reliability as
open work, consistent with this project's stated policy of reporting
measured results only at the scale actually validated.
\end{abstract}

\begin{IEEEkeywords}
Mixture-of-Experts, distributed training, fault tolerance, GPU kernels,
Triton, elastic training, expert parallelism, configuration validation
\end{IEEEkeywords}

\section{Introduction}

Training large Mixture-of-Experts (MoE) models at hyperscale --- tens of
thousands of GPUs --- exposes systems challenges beyond raw throughput:
continuous node failure, silent collective and sharding bugs, misleading
utilization metrics under sparse activation, and the need for seamless
recovery with expert re-ownership \cite{shoeybi2019megatron,
rasley2020deepspeed}. moe-engine is a runtime built specifically around the
premise that at hyperscale, node death is not an exception to be handled
but a continuous background condition the system must remain correct
under.

This paper is the third report on the system. The first version described
the initial Triton router kernel, 4D mesh foundations, and elastic recovery
mechanisms. The second reported production-style multi-process pipeline
parallelism with a full 1F1B schedule, but explicitly declined to present
GPU throughput data, stating that sustained cluster access was not yet
available and that the authors preferred to report no numbers rather than
illustrative ones. This paper reports the closing of that specific gap at
single-GPU scale, together with a substantial internal refactoring
undertaken in parallel.

Our contributions are:
\begin{itemize}
\item The first real-GPU validation of the system's Triton router kernel,
  which surfaced a compilation defect (Section~\ref{sec:kernel}) invisible
  to any CPU-only test, together with measured throughput data
  (Section~\ref{sec:eval}).
\item A structural decomposition of a 1{,}165-line distributed-systems
  module into eight single-responsibility components with full backward
  compatibility (Section~\ref{sec:arch}).
\item An expert-capacity enforcement mechanism following Switch
  Transformer/GShard semantics \cite{fedus2021switch, lepikhin2020gshard},
  implemented via a vectorized group-cumulative-count primitive
  (Section~\ref{sec:capacity}).
\item A root-caused case study of a silent configuration-validation
  failure with a generalizable lesson for systems pairing optional
  dependencies with safety-critical logic (Section~\ref{sec:config}).
\item An expanded test suite (148$\to$348 tests) incorporating
  property-based testing and mocked collective communication
  (Section~\ref{sec:testing}).
\end{itemize}

\section{Related Work}

Megatron-LM \cite{shoeybi2019megatron} and DeepSpeed \cite{rasley2020deepspeed}
established tensor- and pipeline-parallel training at scale for dense
models. Switch Transformer \cite{fedus2021switch} and GShard
\cite{lepikhin2020gshard} introduced sparse expert routing with capacity-based
token dropping, which we adopt directly for our capacity enforcement
mechanism (Section~\ref{sec:capacity}). Tutel \cite{hwang2023tutel}
addressed adaptive expert parallelism at scale. Prior systems in this space
have generally reported results at the multi-node scale their
infrastructure was designed for; our project differs primarily in scope
and maturity level --- it targets research-grade reproducibility on
resources ranging from a single CPU core to, eventually, a full cluster,
and explicitly reports which scale each result was obtained at rather than
presenting cluster-scale projections without cluster-scale data.

\section{System Architecture}
\label{sec:arch}

The system's high-level control flow --- training loop, distributed MoE
layer, elastic training harness, telemetry --- is unchanged from prior
versions. The principal architectural change in this work is a
decomposition of the distributed-systems package, previously a single
1{,}165-line file containing every parallelism primitive, into eight
modules (Table~\ref{tab:modules}), each with one responsibility. A
68-line backward-compatibility shim re-exports every public symbol from
the original module path, so existing call sites require no modification.

\begin{table}[t]
\centering
\caption{Distributed module decomposition}
\label{tab:modules}
\begin{tabular}{@{}p{0.32\linewidth}p{0.58\linewidth}@{}}
\toprule
\textbf{Module} & \textbf{Responsibility} \\
\midrule
\texttt{mesh.py} & Topology, device mesh, comm.\ streams \\
\texttt{tensor\_parallel.py} & Column/Row-parallel linear layers \\
\texttt{sequence\_parallel.py} & SP scatter/gather, fused all-gather \\
\texttt{expert\_parallel.py} & All-to-all dispatch/combine \\
\texttt{pipeline\_parallel.py} & 1F1B pipeline schedule \\
\texttt{data\_parallel.py} & FSDP2 application \\
\texttt{router.py} & High-level router interface \\
\texttt{moe\_layer.py} & Layer orchestration \\
\bottomrule
\end{tabular}
\end{table}

Model definitions were additionally extracted into a separate package with
a decorator-based registry (\texttt{register\_model},
\texttt{build\_model\_from\_config}), following the pattern used by
Megatron-LM and torchTitan to decouple training infrastructure from model
research code: the training entrypoint dispatches on a configuration
string rather than importing a specific model class.

\section{Triton Router Kernel}
\label{sec:kernel}

The router kernel fuses gating matmul, softmax, top-$K$ selection, and
combine-weight renormalization into a single Triton kernel operating in
one HBM pass, unchanged architecturally since our first report. This
section describes a defect surfaced only when the kernel was, for the
first time, compiled and executed on real GPU hardware as part of this
work.

\subsection{The constexpr Defect}

Both the forward and backward kernels unroll their top-$K$ selection loop
using \texttt{tl.static\_range}, which requires the loop bound to be a
compile-time \texttt{constexpr} value. The kernel signatures declared $K$
as an ordinary runtime integer. On the CPU-only reference implementation
used throughout prior continuous integration --- plain PyTorch, never
invoking the Triton compiler --- this had no observable effect. The first
real-GPU compilation attempt failed with a Triton
\texttt{CompilationError}.

We highlight this defect because of its generalizable failure mode:
\emph{a CPU reference fallback for hardware-specific code is exactly what
allows a hardware-only defect to survive an arbitrary number of CI runs},
because CI never exercises the code path where the defect resides. The
fix --- declaring $K$ as \texttt{constexpr} in both kernel signatures ---
is minor; the broader lesson motivated two regression tests: a
CPU-executable static assertion that both kernel signatures declare $K$ as
\texttt{constexpr}, and a GPU-gated test that compiles and runs the kernel
when CUDA is available.

\section{4D Parallelism}

Data, expert, tensor, and pipeline parallelism compose via a device mesh
constructed once at startup. Tensor parallelism shards SwiGLU expert
weights following the standard column/column/row convention (both gate and
up projections column-parallel, down projection row-parallel), unchanged
from prior work. Pipeline parallelism uses a three-phase 1F1B schedule
(warmup, steady-state, drain) with real \texttt{torch.distributed.send}/
\texttt{recv} on a dedicated process group and explicit micro-batch
tagging, validated by 2-rank \texttt{mp.spawn} tests; this component is
unchanged in the present work and has not yet been exercised under
simulated node failure.

\section{Expert Capacity Management}
\label{sec:capacity}

The principal functional addition in this work is an optional expert
capacity enforcement mechanism following Switch Transformer/GShard
semantics: each expert accepts at most
\begin{equation}
c_e = \left\lceil \gamma \cdot \frac{N K}{E} \right\rceil
\end{equation}
tokens, where $\gamma$ is a configurable capacity factor, $N$ is the token
count, $K$ is active experts per token, and $E$ is total expert count.
Tokens beyond this budget receive zero combine weight for the
corresponding top-$K$ slot rather than being processed.

\begin{algorithm}[t]
\caption{\texttt{cumcount}: position of appearance within group}
\label{alg:cumcount}
\begin{algorithmic}[1]
\REQUIRE groups $\in \mathbb{Z}^N$ (expert assignment per token)
\ENSURE position $\in \mathbb{Z}^N$ (0-indexed rank within group)
\STATE order $\gets$ \textsc{StableArgsort}(groups)
\STATE sorted $\gets$ groups[order]
\STATE isStart$[0] \gets$ true; isStart$[i] \gets$ sorted$[i] \ne$ sorted$[i{-}1]$
\STATE startIdx $\gets$ \textsc{Where}(isStart, $i$, 0)
\STATE startIdx $\gets$ \textsc{CumMax}(startIdx)
\STATE posSorted $\gets i - $ startIdx
\STATE position[order] $\gets$ posSorted
\RETURN position
\end{algorithmic}
\end{algorithm}

Algorithm~\ref{alg:cumcount} gives the core primitive, \texttt{cumcount},
which computes each token's zero-indexed position of appearance within its
assigned expert's queue using a stable sort followed by a cumulative
maximum, avoiding both an explicit per-expert loop and a dependency on
\texttt{scatter\_reduce} variants whose availability differs across
PyTorch releases. The capacity drop mask is then
$\text{drop}[n,k] = \texttt{cumcount}(\text{idx}[:,k])[n] \ge c_e$,
applied independently per top-$K$ slot. The mechanism is disabled by
default; enabling it changes no existing behavior for callers who do not
opt in. Twenty-two dedicated tests cover the primitive, mask construction,
and full-layer integration including gradient flow through non-dropped
slots.

A companion configuration exercising $E{=}256$, $K{=}8$ --- a fine-grained
expert regime four times the expert count of our existing production
configuration --- validates the mechanism's correctness at reduced
dimensions; its throughput and memory behavior at full scale on real
hardware has not been measured (Section~\ref{sec:limitations}).

\section{Configuration Validation}
\label{sec:config}

\subsection{Schema Migration}

Prior configuration handling used an unvalidated flat dictionary: an
invalid value such as $\text{top\_k} > \text{num\_experts}$ would not
surface until the router kernel failed with an opaque index error,
potentially minutes into training. We replaced this with a statically
validated hierarchical schema with cross-field validators enforcing, among
other constraints, $\text{top\_k} \le \text{num\_experts}$ and
$\text{warmup\_steps} < \text{max\_steps}$, reporting the offending field
path and constraint description at load time.

\subsection{A Silent-Degradation Case Study}

We report the following defect in detail for its methodological interest.
The configuration module included a fallback: if the validation library
were unavailable, the module would degrade to an unvalidated pass-through
rather than fail to import, intended to preserve importability in minimal
environments. The dependency, however, was never declared in the package's
runtime requirements. Continuous integration installed exactly what was
declared, so the validation library was never present, and every
configuration --- including deliberately invalid ones constructed by the
test suite specifically to verify that validation errors fire --- silently
passed through the unvalidated path. Eight assertions that a specific
exception type would be raised instead observed no exception, for an
entire release cycle, without any test reporting failure, because the
tests never exercised the code path they were written to test.

\begin{table}[t]
\centering
\caption{Silent-degradation defects and their fixes}
\label{tab:bugs}
\small
\begin{tabular}{@{}p{0.46\linewidth}p{0.46\linewidth}@{}}
\toprule
\textbf{Root cause} & \textbf{Fix} \\
\midrule
Validation library never declared as a runtime dependency; optional-import fallback silently disabled all validation & Declared as a hard dependency; fallback path deleted entirely --- module now fails loudly on import if absent \\
\addlinespace
\texttt{yaml.safe\_load("1e-5")} returns a string under YAML 1.1 grammar, not a float & Native \texttt{int}/\texttt{float} coercion attempted before YAML fallback \\
\bottomrule
\end{tabular}
\end{table}

The general lesson (Table~\ref{tab:bugs}) we draw is that \emph{a component
whose stated purpose is catching invalid input early should not be able to
silently stop performing that function}; an import error is a safer
failure mode than a silently permissive one. We deleted the fallback path
entirely rather than fixing the dependency declaration alone.

\section{Fault Tolerance and Checkpointing}

Asynchronous two-tier checkpointing (pinned host memory, NVMe with
\texttt{O\_DIRECT} and atomic rename, S3) and round-robin expert resharding
on node loss are unchanged in core algorithm from prior work. We added
explicit schema versioning to checkpoint metadata, embedding the runtime's
package and PyTorch versions (read from their respective single sources of
truth, not hardcoded) alongside a monotonically increasing schema version
integer, with graceful field defaulting on version mismatch in either
direction. The storage-stall chaos scenario remains fully verified and, as
reported in Section~\ref{sec:eval}, now has real-GPU confirmation. The
node-kill-and-recovery scenario is unchanged at approximately 85\% pass
rate due to a Gloo process-group re-initialization race, whose fix (NCCL
migration in the chaos harness) requires GPU hardware not available for
this cycle.

\section{Testing Infrastructure}
\label{sec:testing}

The CPU-only test tier grew from 148 to 348 passing tests over this arc.
Three additions specifically target correctness properties that do not
require GPU or multi-process resources: (i) nine Hypothesis-based
property tests verifying token conservation, weight normalization,
complete and non-overlapping expert ownership under arbitrary
$(E, \text{ep\_size})$, and configuration invariants, each checked against
fifty randomly generated inputs rather than fixed examples; (ii) a
thread-based mock of \texttt{all\_to\_all\_single} enabling
expert-parallel dispatch logic to be exercised at $\text{ep\_size}{>}1$
within a single process, complementary to (not a replacement for) existing
2-rank \texttt{mp.spawn} ground-truth tests; and (iii) dedicated test
suites for the high-level router interface and model registry (33 and 20
tests respectively), following the principle that each public interface
boundary warrants tests that would fail if that boundary's contract were
violated, independent of underlying-implementation coverage.

\section{Evaluation}
\label{sec:eval}

\subsection{Real-GPU Validation}

We report results from a single NVIDIA T4 (16~GB HBM2, 65 TFLOPS FP32),
the first sustained GPU validation for this system. Token conservation
held with zero violations across 100 seeds on CUDA, matching the CPU
reference. The storage-stall chaos scenario achieved a 10-of-10 (100\%)
pass rate under real-GPU execution.

\begin{table}[t]
\centering
\caption{Router forward throughput: T4 GPU vs.\ CPU reference}
\label{tab:speedup}
\small
\begin{tabular}{@{}rrrrrrr@{}}
\toprule
$N$ & $H$ & $E$ & $K$ & CPU (M/s) & GPU (M/s) & Speedup \\
\midrule
512  & 256  & 16 & 2 & 0.747 & 2.141 & 2.9$\times$ \\
1024 & 512  & 32 & 2 & 0.421 & 3.644 & 8.7$\times$ \\
2048 & 1024 & 64 & 2 & 0.236 & 4.832 & 20.4$\times$ \\
4096 & 2048 & 64 & 4 & 0.056 & 4.454 & \textbf{80.1$\times$} \\
\bottomrule
\end{tabular}
\end{table}

Table~\ref{tab:speedup} reports forward-only router throughput at four
configurations. Speedup increases super-linearly with problem size, from
2.9$\times$ at $N{=}512$ to 80.1$\times$ at $N{=}4096, H{=}2048$ --- the
configuration closest to production MoE hidden dimensions. We attribute
this to the single-HBM-pass kernel's bandwidth utilization improving with
working-set size, unmatched by the single-threaded, unvectorized CPU
reference path.

We emphasize scope: these results confirm kernel correctness and
CPU/GPU scaling on real hardware for the first time, but do not constitute
production-scale MFU. The T4 is inference-oriented, with roughly
$15\times$ lower BF16 peak throughput than an H100 SXM5, and measured MFU
at smoke-test batch sizes is on the order of $10^{-6}$ --- expected and
uninformative at this scale, reported for completeness rather than as a
substantive result.

\subsection{MoE Layer Overhead}

Full distributed-layer forward latency, measured against a dense
single-expert baseline at matched $(B,S,H,F)$, showed 21.8$\times$ to
26.1$\times$ overhead on the T4 and 9.1$\times$ to 16.2$\times$ on CPU at
small batch sizes ($B{\le}4$). This range is dominated by kernel-launch
latency at these sizes and by a parameter-count mismatch between the
dense baseline and the MoE configuration compared against it; it should
be read as a latency-floor measurement rather than a throughput-per-parameter
comparison.

\section{Limitations and Future Work}
\label{sec:limitations}

Consistent with this project's reporting policy, we state open items
explicitly rather than eliding them. \emph{Multi-node GPU data} does not
yet exist beyond the single-T4 results above; this remains the top
priority for future work and requires sustained cluster access.
\emph{Chaos Scenario A} (node kill and recovery) remains at approximately
85\% due to the unresolved Gloo race noted in Section~VII; its fix
requires GPU hardware to validate. The \emph{expert capacity mechanism} of
Section~\ref{sec:capacity} is validated for correctness only at reduced
scale; its behavior at $E{=}256$ under real dispatch bandwidth constraints
is unmeasured, as is a next-best-expert re-routing policy as an
alternative to outright dropping. \emph{Pipeline parallelism under
simulated node failure} has not been exercised. \emph{Non-divisible
sequence lengths} under sequence parallelism remain unsupported.

\section{Conclusion}

This paper reports closing a specific, previously stated gap: real-GPU
validation of moe-engine's router kernel and chaos-recovery mechanism,
with reproducible single-device measurements rather than illustrative
estimates. That validation surfaced a genuine defect invisible to any
CPU-only test, illustrating a general risk in systems that maintain
hardware-specific code behind a CPU fallback. In parallel work, a
structural refactoring surfaced a second, independent defect --- a
configuration validator silently non-functional for a full release cycle
due to an undeclared dependency --- which we argue generalizes beyond this
system to any component pairing safety-critical validation with an
optional-dependency fallback. The system's multi-node evaluation gap
remains open for the reason it has been open throughout this project's
reporting history: access to sustained cluster resources, not algorithmic
readiness.

\section*{Acknowledgment}
The author thanks the PyTorch, Triton, and distributed systems
communities, the maintainers of Pydantic and Hypothesis, and researchers
who have openly shared large-scale training failure post-mortems.

\begin{thebibliography}{9}

\bibitem{shoeybi2019megatron}
M. Shoeybi, M. Patwary, R. Puri, P. LeGresley, J. Casper, and B.
Catanzaro, ``Megatron-LM: Training multi-billion parameter language
models using model parallelism,'' \emph{arXiv:1909.08053}, 2019.

\bibitem{rasley2020deepspeed}
J. Rasley, S. Rajbhandari, O. Ruwase, and Y. He, ``DeepSpeed: System
optimizations enable training deep learning models with over 100 billion
parameters,'' in \emph{Proc. 26th ACM SIGKDD Conf.}, 2020, pp. 3505--3506.

\bibitem{fedus2021switch}
W. Fedus, B. Zoph, and N. Shazeer, ``Switch transformers: Scaling to
trillion parameter models with simple and efficient sparsity,''
\emph{arXiv:2101.03961}, 2021.

\bibitem{lepikhin2020gshard}
D. Lepikhin, H. Lee, Y. Xu, D. Chen, O. Firat, Y. Huang, M. Krikun, N.
Shazeer, and Z. Chen, ``GShard: Scaling giant models with conditional
computation and automatic sharding,'' \emph{arXiv:2006.16668}, 2020.

\bibitem{hwang2023tutel}
C. Hwang \emph{et al.}, ``Tutel: Adaptive mixture-of-experts at scale,''
\emph{arXiv:2206.03382}, 2023.

\end{thebibliography}

\end{document}
